\chapter{Introduction}
\label{ch:introduction}

\section{Purpose and Scope}

The SIRC-1 documentation is split into three manuals, each with its own scope:

\begin{table}[H]
	\centering
	\begin{tabular}{lp{9cm}}
		\toprule
		\textbf{Manual}                    & \textbf{Scope}                                                                                \\
		\midrule
		SIRC Architecture Reference Manual & ISA, instruction encodings, canonical assembly notation, exceptions, programmer-visible state \\
		SIRC Assembly Programmer's Manual  & Complete assembly syntax, directives, macros, pseudoinstructions (this manual)                \\
		SIRC Hardware Reference Manual     & Pinout, bus interface, timing, electrical characteristics                                     \\
		\bottomrule
	\end{tabular}
	\caption{The SIRC manual set}
	\label{tab:manual-set}
\end{table}

This manual covers the language and the toolchain. It does not define what an instruction does, which registers exist,
or how exceptions work. For those, see the \archmanual{}; this manual refers to it by chapter name.

\section{The Toolchain}

A SIRC program is built in three steps, each with its own tool:

\begin{enumerate}
	\item The \textbf{assembler} reads one assembly source file and writes an object file. It encodes every instruction and data
	      directive, and records every label and every reference to a label.
	\item The \textbf{linker} reads one or more object files, joins them end to end, resolves every reference to a label, and
	      writes a flat binary program plus a debug map.
	\item The \textbf{simulator} loads the binary at address zero, maps whatever memory the program needs, and runs it.
\end{enumerate}

\begin{table}[H]
	\centering
	\begin{tabular}{lll}
		\toprule
		\textbf{Tool}      & \textbf{Reads}           & \textbf{Writes}                                          \\
		\midrule
		\texttt{assembler} & \texttt{.sasm} source    & Object file (\texttt{.o})                                \\
		\texttt{linker}    & One or more object files & Binary (\texttt{.bin}) and debug map (\texttt{.bin.dbg}) \\
		\texttt{sirc\_vm}  & Binary                   & Register dump, terminal output                           \\
		\bottomrule
	\end{tabular}
	\caption{The SIRC toolchain}
	\label{tab:toolchain}
\end{table}

Chapter~\ref{ch:toolchain} describes how to run each tool.

\section{Conventions}

Source listings are shown in typewriter font. Mnemonics are bold typewriter, as in the architecture manual. File
extensions are conventions, not requirements: the tools do not look at them.
